\documentclass[12pt]{article}

\usepackage[a4paper, top=2cm, bottom=2cm, left=2.5cm, right=2.5cm]{geometry}

\usepackage{amsmath, amssymb}
\usepackage{tikz}
\usetikzlibrary{arrows.meta,positioning,fit,calc}
\usepackage{multirow}
\usepackage{array}
\usepackage{siunitx}
\usepackage{enumitem}
\usepackage{float}
\usepackage{fontspec}
\usetikzlibrary{decorations.pathmorphing}
\usepackage{pdflscape}
\usepackage{tabularx}
\usepackage{booktabs}
\usepackage{graphicx}
\usepackage{listings}
\usepackage[table]{xcolor}
\usepackage{hyperref}
\usepackage{bookmark}
\usepackage{pdfpages}

% ---------------- FONTS ----------------
% Body font
\setmainfont[
    UprightFont = Handlee-Regular.ttf,
    AutoFakeBold = 2.5
]{Handlee}

% Heading font
\newfontfamily\headingfont{PT.ttf}

% verbatim font
\setmonofont{DejaVu Sans Mono} % or Noto Sans Mono

% ---------------- HANDWRITTEN MATH ----------------
\usepackage{mathastext}
\Mathastext
\everymath{\displaystyle}

% Optional spacing tweaks (more natural look)
% \setlength{\thinmuskip}{2mu}
% \setlength{\medmuskip}{2mu}

% ---------------- GENERAL ----------------
\setlength{\parindent}{0pt}
\pagenumbering{gobble}
\linespread{1.2}


\newcommand{\mychapter}[2]{%
    \phantomsection
    \hypertarget{#2}{}%
    \bookmark[level=0,dest=#2]{#1}%
    \begin{center}
        {\headingfont\Huge #1}
        \par
        \vspace{4pt}
        \hrule
    \end{center}
}

\newcommand{\mysection}[3]{%
    \phantomsection
    \hypertarget{#3}{}%
    \bookmark[level=1,dest=#3]{#1\space #2}%
    \newpage
    \noindent{\headingfont\Large #1\space #2\par}
    \vspace{4pt}
    \hrule
    \vspace{15pt}
}

\newcommand{\mysubsection}[2]{%
    \vspace{5px}
    \noindent{\headingfont\large #1\space #2\par}
    \vspace{4pt}
    \hrule
    \vspace{15pt}
}

\newcommand{\insertfigure}[2]{%
    \begin{figure}[H]
        \centering
        \fbox{\includegraphics[width=#1\textwidth]{../2. Figures/#2}}
    \end{figure}
}

\begin{document}
\vspace*{\fill}

\insertfigure{0.6}{cover.png}


\vspace{1cm}

\mychapter{Chapter 3 : Properties of Pure Substance}{5f4y5fdh}

This chapter introduces the concept of pure substances and examines their different phases and phase-change processes. It develops the use of property diagrams, including $T$-$v$, $P$-$v$, and $P$-$T$ diagrams, to understand thermodynamic behavior. The chapter also explains property tables and the determination of thermodynamic properties of pure substances. Finally, it introduces the ideal-gas equation of state, compressibility factor, and common equations of state.

\vspace*{\fill}

\newpage

\begin{center}
    {\headingfont\Huge Chapter Tree}
    \par
    \vspace{4pt}
    \hrule
\end{center}
\begin{verbatim}
│
├── 3–1 PURE SUBSTANCE
│   Fixed chemical composition throughout the substance.
│
├── 3–2 PHASES OF A PURE SUBSTANCE
│   Describes solid, liquid, and gas phases microscopically.
│
├── 3–3 PHASE-CHANGE PROCESSES OF PURE SUBSTANCES
│   Explains phase changes using water as the model substance.
│
├── 3–4 PROPERTY DIAGRAMS FOR PHASE-CHANGE PROCESSES
│   Uses property diagrams to represent phase-change behavior.
│
├── 3–5 PROPERTY TABLES
│   Presents thermodynamic properties through organized tabulated data.
│
├── 3–6 THE IDEAL-GAS EQUATION OF STATE
│   Introduces the ideal-gas relation and its practical applicability.
│
├── 3–7 COMPRESSIBILITY FACTOR—A MEASURE OF DEVIATION FROM IDEAL
│        -GAS BEHAVIOR
│   Quantifies deviations of real gases from ideal behavior.
│
└── 3–8 OTHER EQUATIONS OF STATE
    Introduces more accurate real-gas equations of state.

    
\end{verbatim}















\mysection{3.1}{Pure Substance}{3.1}

A \textbf{pure substance} is a substance that has a fixed chemical composition throughout.

\textbf{\\ Examples}

Water, nitrogen, helium, and carbon dioxide are pure substances.

\textbf{\\ Homogeneous Mixtures}

A pure substance does not necessarily have to consist of a single chemical element or compound.

A mixture of different chemical elements or compounds is also considered a pure substance if the mixture is \textbf{homogeneous}, meaning that its chemical composition is uniform throughout.

For example, air is a mixture of several gases, but it can be considered a pure substance because its composition is uniform throughout.

\textbf{\\ Nonhomogeneous Mixtures}

A mixture is not a pure substance when its composition is not uniform throughout.

For example, a mixture of oil and water is not a pure substance because oil is not soluble in water. The oil collects above the water, producing two chemically dissimilar regions.

\textbf{\\ Pure Substance with Multiple Phases}

A mixture containing two or more phases can still be a pure substance if the \textbf{chemical composition of all phases is the same}.

\begin{center}
    \begin{tabular}{|c|c|}
        \hline
        \textbf{Mixture} & \textbf{Pure Substance?} \\
        \hline
        Ice + liquid water & Yes \\
        \hline
        Liquid air + gaseous air & No \\
        \hline
        Oil + water & No \\
        \hline
    \end{tabular}
\end{center}

\textbf{\\ Ice--Water Mixture}

Ice and liquid water form a pure substance because both phases have the same chemical composition, namely $\mathrm{H_2O}$.

\textbf{\\ Liquid--Gaseous Air Mixture}

A mixture of liquid air and gaseous air is not a pure substance because the composition of liquid air differs from that of gaseous air.

This difference occurs because the different components of air condense at different temperatures at a specified pressure.

\insertfigure{0.6}{Figure : 1.png}



\textbf{\\ Key Criterion}

The essential requirement for a pure substance is:

\begin{equation*}
\boxed{\text{Uniform chemical composition throughout}}
\end{equation*}

A substance may contain:
\begin{itemize}
    \item A single chemical element,
    \item A single chemical compound,
    \item A homogeneous mixture of several components, or
    \item Multiple phases, provided all phases have the same chemical composition.
\end{itemize}

\textbf{\\ Important Distinction}

The number of components or phases alone does not determine whether a substance is pure.

\begin{equation*}
\boxed{\text{Same composition in all regions/phases} \;\Rightarrow\; \text{Pure substance}}
\end{equation*}

\begin{equation*}
\boxed{\text{Different composition in different regions/phases} \;\Rightarrow\; \text{Not a pure substance}}
\end{equation*}





















\mysection{3.2}{Phases of a Pure Substance}{3.2}

Substances exist in different phases depending on the prevailing conditions of temperature and pressure. The three principal phases are \textbf{solid}, \textbf{liquid}, and \textbf{gas}.

\textbf{\\ Examples of Phases}

At room temperature and pressure:
\begin{center}
    \begin{tabular}{|c|c|}
        \hline
        \textbf{Substance} & \textbf{Phase} \\
        \hline
        Copper & Solid \\
        \hline
        Mercury & Liquid \\
        \hline
        Nitrogen & Gas \\
        \hline
    \end{tabular}
\end{center}

Under different conditions, the same substance may exist in a different phase.

\textbf{\\ Multiple Phases within a Principal Phase}

A substance may have several phases within the same principal phase. These phases can have different molecular structures.

\begin{center}
    \begin{tabular}{|c|c|c|}
        \hline
        \textbf{Substance} & \textbf{Principal Phase} & \textbf{Different Phases} \\
        \hline
        Carbon & Solid & Graphite, Diamond \\
        \hline
        Helium & Liquid & Two liquid phases \\
        \hline
        Iron & Solid & Three solid phases \\
        \hline
        Ice & Solid & Seven phases at high pressures \\
        \hline
    \end{tabular}
\end{center}

\textbf{\\ Definition of a Phase}

A \textbf{phase} is a state of a substance having:
\begin{itemize}
    \item a distinct molecular arrangement,
    \item a homogeneous structure throughout, and
    \item clearly identifiable boundary surfaces separating it from other phases.
\end{itemize}

The two phases of $\mathrm{H_2O}$ in iced water provide an example of distinct phases of the same substance.
\textbf{\\ Comparison of Molecular Arrangements}

\begin{center}
    \begin{tabular}{|c|c|c|c|}
        \hline
        \textbf{Property} & \textbf{Solid} & \textbf{Liquid} & \textbf{Gas} \\
        \hline
        Molecular arrangement & Ordered lattice & Less ordered & No molecular order \\
        \hline
        Molecular position & Relatively fixed & Not fixed & Widely separated \\
        \hline
        Molecular motion & Oscillation & Rotation and translation & Random motion \\
        \hline
        Intermolecular force & Strongest & Intermediate & Weakest \\
        \hline
        Molecular energy & Lowest & Intermediate & Highest \\
        \hline
    \end{tabular}
\end{center}

\textbf{\\ Phase Transformation and Temperature}

The temperature of a substance affects the velocity and momentum of its molecules.

\begin{equation*}
\text{Temperature} \uparrow
\quad \Rightarrow \quad
\text{Molecular velocity} \uparrow
\quad \Rightarrow \quad
\text{Intermolecular forces can be overcome}
\end{equation*}

At sufficiently high temperature, groups of molecules in a solid can escape their fixed positions, initiating melting.

Conversely, a gas must release a large amount of energy to move toward the lower-energy liquid or solid phases.

\insertfigure{1}{Figure : 2}




















\mysection{3.3}{Phase-Change Processes of Pure Substances}{3.3}

There are many practical situations in which two phases of a pure substance coexist in equilibrium. Examples include:
\begin{itemize}
    \item liquid--vapor water in the boiler and condenser of a steam power plant,
    \item liquid--vapor refrigerant in a refrigerator, and
    \item freezing of water in underground pipes.
\end{itemize}

Water is used to demonstrate the basic principles of phase-change processes, but all pure substances exhibit the same general behavior.

\textbf{\\ Compressed Liquid and Saturated Liquid}

Consider a piston--cylinder device containing liquid water at $20^\circ\mathrm{C}$ and $1~\mathrm{atm}$.

At this state, water exists as a liquid and is called a \textbf{compressed liquid}, or \textbf{subcooled liquid}. It is not about to vaporize.

When heat is transferred to the water:
\begin{itemize}
    \item the temperature increases,
    \item the liquid expands slightly,
    \item the specific volume increases, and
    \item the piston moves upward slightly.
\end{itemize}

As more heat is transferred, the temperature eventually reaches $100^\circ\mathrm{C}$. At this point, water is still entirely liquid, but any additional heat causes some of the liquid to vaporize.

A liquid that is \textbf{about to vaporize} is called a \textbf{saturated liquid}.

Therefore:
\begin{equation*}
\boxed{\text{Compressed liquid} \rightarrow \text{Saturated liquid} \rightarrow \text{Vaporization}}
\end{equation*}

For water at $1~\mathrm{atm}$:
\begin{equation*}
T_{\mathrm{sat}} = 100^\circ\mathrm{C}
\end{equation*}

\textbf{\\ Saturated Liquid--Vapor Mixture}

Once boiling begins at constant pressure, the temperature remains constant until the liquid has been completely vaporized.

During this phase-change process:
\begin{itemize}
    \item heat continues to be transferred to the substance,
    \item the temperature remains constant,
    \item the volume increases substantially, and
    \item the amount of liquid decreases while the amount of vapor increases.
\end{itemize}

At an intermediate state, both liquid and vapor exist simultaneously. This is called a \textbf{saturated liquid--vapor mixture}.

The liquid and vapor phases coexist in equilibrium.


\textbf{\\ Saturated Vapor}

Continued heat transfer eventually vaporizes the last drop of liquid.

At this point, the cylinder contains only vapor, but the vapor is on the \textbf{borderline of the liquid phase}.

Any heat loss from this vapor will cause some of it to condense.

A vapor that is \textbf{about to condense} is called a \textbf{saturated vapor}.

Therefore:
\begin{equation*}
\boxed{\text{Saturated liquid} \rightarrow
\text{Saturated liquid--vapor mixture} \rightarrow
\text{Saturated vapor}}
\end{equation*}

The states between saturated liquid and saturated vapor constitute the saturated liquid--vapor mixture region.

\textbf{\\ Superheated Vapor}

After the last drop of liquid has vaporized, further heat transfer causes both:
\begin{itemize}
    \item temperature to increase, and
    \item specific volume to increase.
\end{itemize}

A vapor that is \textbf{not about to condense} is called a \textbf{superheated vapor}.

For example, at $1~\mathrm{atm}$, water at $300^\circ\mathrm{C}$ is a superheated vapor.

As long as the temperature remains above the saturation temperature at the specified pressure, cooling the vapor does not immediately cause condensation.

\textbf{\\ Constant-Pressure Heating Process}

For water heated at constant pressure, the complete process can be represented as:


\begin{equation*}
\hspace*{-2cm}
\boxed{
\text{Compressed liquid}
\rightarrow
\text{Saturated liquid}
\rightarrow
\text{Saturated liquid--vapor mixture}
\rightarrow
\text{Saturated vapor}
\rightarrow
\text{Superheated vapor}
}
\end{equation*}

For water at $1~\mathrm{atm}$:

\begin{center}
    \begin{tabular}{|c|c|c|}
        \hline
        \textbf{State} & \textbf{Condition} & \textbf{Temperature} \\
        \hline
        1 & Compressed liquid & $20^\circ\mathrm{C}$ \\
        \hline
        2 & Saturated liquid & $100^\circ\mathrm{C}$ \\
        \hline
        3 & Saturated liquid--vapor mixture & $100^\circ\mathrm{C}$ \\
        \hline
        4 & Saturated vapor & $100^\circ\mathrm{C}$ \\
        \hline
        5 & Superheated vapor & $300^\circ\mathrm{C}$ \\
        \hline
    \end{tabular}
\end{center}

During the phase-change process from state 2 to state 4:

\begin{equation*}
\boxed{P=\text{constant},\qquad T=\text{constant}}
\end{equation*}

\insertfigure{0.6}{Figure : 3.png}

\textbf{\\ Reversing the Process}

If the heating process is reversed by cooling the water while maintaining the same pressure, the water retraces the same path back to its initial state.

The amount of heat released during the reverse process exactly matches the amount of heat added during the heating process.

\textbf{\\ Water and Steam in Thermodynamics}

In everyday usage, water generally refers to liquid water and steam refers to water vapor.

In thermodynamics, both \textbf{water} and \textbf{steam} generally refer to the same chemical substance:

\begin{equation*}
\boxed{\mathrm{H_2O}}
\end{equation*}

\textbf{\\ Saturation Temperature and Saturation Pressure}

The statement ``water boils at $100^\circ\mathrm{C}$'' is incomplete.

The correct statement is:

\begin{equation*}
\boxed{\text{Water boils at }100^\circ\mathrm{C}\text{ at }1~\mathrm{atm}}
\end{equation*}

The boiling temperature depends on pressure.

For example:
\begin{center}
    \begin{tabular}{|c|c|}
        \hline
        \textbf{Pressure} & \textbf{Boiling Temperature of Water} \\
        \hline
        $1~\mathrm{atm}$ & $100^\circ\mathrm{C}$ \\
        \hline
        $500~\mathrm{kPa}$ & $151.8^\circ\mathrm{C}$ \\
        \hline
    \end{tabular}
\end{center}

At a given pressure, the temperature at which a pure substance changes phase is called the \textbf{saturation temperature}, $T_{\mathrm{sat}}$.

At a given temperature, the pressure at which a pure substance changes phase is called the \textbf{saturation pressure}, $P_{\mathrm{sat}}$.

\begin{equation*}
\boxed{T_{\mathrm{sat}}=\text{temperature at which phase change occurs at a specified pressure}}
\end{equation*}

\begin{equation*}
\boxed{P_{\mathrm{sat}}=\text{pressure at which phase change occurs at a specified temperature}}
\end{equation*}

For water:

\begin{equation*}
P_{\mathrm{sat}}=101.325~\mathrm{kPa}
\quad\Longleftrightarrow\quad
T_{\mathrm{sat}}=99.97^\circ\mathrm{C}
\end{equation*}

At $100.00^\circ\mathrm{C}$, the saturation pressure is approximately $101.42~\mathrm{kPa}$ according to the ITS-90 scale.

\textbf{\\ Saturation Tables}

Saturation tables provide the relationship between saturation pressure and saturation temperature.

For water, selected values are:

\begin{center}
    \begin{tabular}{|c|c|}
        \hline
        \textbf{Temperature, $T$ ($^\circ\mathrm{C}$)} & \textbf{Saturation Pressure, $P_{\mathrm{sat}}$ (kPa)} \\
        \hline
        $-10$ & $0.260$ \\
        \hline
        $0$ & $0.611$ \\
        \hline
        $10$ & $1.23$ \\
        \hline
        $20$ & $2.34$ \\
        \hline
        $25$ & $3.17$ \\
        \hline
        $50$ & $12.35$ \\
        \hline
        $100$ & $101.3$ \\
        \hline
        $150$ & $475.8$ \\
        \hline
        $200$ & $1554$ \\
        \hline
        $250$ & $3973$ \\
        \hline
        $300$ & $8581$ \\
        \hline
    \end{tabular}
\end{center}

Thus, at $25^\circ\mathrm{C}$, water changes phase at approximately:

\begin{equation*}
P_{\mathrm{sat}}=3.17~\mathrm{kPa}
\end{equation*}

Water can be frozen by reducing its pressure below approximately:

\begin{equation*}
P=0.61~\mathrm{kPa}
\end{equation*}

At $250^\circ\mathrm{C}$, water boils when its pressure is maintained at approximately:

\begin{equation*}
P_{\mathrm{sat}}\approx3973~\mathrm{kPa}
\end{equation*}

\textbf{\\ Latent Heat}

A large amount of energy is required to melt a solid or vaporize a liquid.

The energy absorbed or released during a phase-change process is called the \textbf{latent heat}.

\textbf{\\ Latent Heat of Fusion}

The energy absorbed during melting is called the \textbf{latent heat of fusion}.

It is equal in magnitude to the energy released during freezing.

\begin{equation*}
\boxed{\text{Melting: energy absorbed}}
\end{equation*}

\begin{equation*}
\boxed{\text{Freezing: energy released}}
\end{equation*}

For water at $1~\mathrm{atm}$:

\begin{equation*}
h_{\mathrm{fusion}}=333.7~\mathrm{kJ/kg}
\end{equation*}

\textbf{\\ Latent Heat of Vaporization}

The energy absorbed during vaporization is called the \textbf{latent heat of vaporization}.

It is equal in magnitude to the energy released during condensation.

\begin{equation*}
\boxed{\text{Vaporization: energy absorbed}}
\end{equation*}

\begin{equation*}
\boxed{\text{Condensation: energy released}}
\end{equation*}

For water at $1~\mathrm{atm}$:

\begin{equation*}
h_{\mathrm{vap}}=2256.5~\mathrm{kJ/kg}
\end{equation*}

The magnitude of latent heat depends on the temperature or pressure at which the phase change occurs.

\textbf{\\ Liquid--Vapor Saturation Curve}

During a phase-change process, pressure and temperature are dependent properties.

Their relationship can be written as:

\begin{equation*}
P_{\mathrm{sat}}=f(T_{\mathrm{sat}})
\end{equation*}

A plot of $P_{\mathrm{sat}}$ versus $T_{\mathrm{sat}}$ is called the \textbf{liquid--vapor saturation curve}.

The saturation curve is characteristic of a pure substance.

\insertfigure{0.4}{Figure : 4.png}

From the saturation curve:

\begin{equation*}
\boxed{P_{\mathrm{sat}}\uparrow \quad\Rightarrow\quad T_{\mathrm{sat}}\uparrow}
\end{equation*}

Therefore, a substance boils at a higher temperature when its pressure is increased.

\textbf{\\ Effect of Pressure on Boiling}

Increasing pressure increases the boiling temperature.

This principle is used in pressure cookers. At approximately $3~\mathrm{atm}$ absolute pressure, water boils at about $134^\circ\mathrm{C}$.

The higher boiling temperature allows food to cook faster and can reduce cooking time and energy consumption.

For example, beef stew may take approximately $1$--$2$ hours in a regular pan at $1~\mathrm{atm}$ but only about $20$ minutes in a pressure cooker operating at $3~\mathrm{atm}$ absolute pressure.

\textbf{\\ Effect of Elevation on Boiling}

Atmospheric pressure decreases as elevation increases.

Since boiling temperature depends on pressure, the boiling temperature of water also decreases with elevation.

At an elevation of $2000~\mathrm{m}$:
\begin{equation*}
P_{\mathrm{atm}}=79.50~\mathrm{kPa}
\end{equation*}

and the corresponding boiling temperature is:

\begin{equation*}
T_{\mathrm{boil}}=93.3^\circ\mathrm{C}
\end{equation*}

At sea level:

\begin{equation*}
P_{\mathrm{atm}}\approx101.33~\mathrm{kPa},
\qquad
T_{\mathrm{boil}}\approx100^\circ\mathrm{C}
\end{equation*}

The boiling temperature drops by a little more than $3^\circ\mathrm{C}$ for every $1000~\mathrm{m}$ increase in elevation.

Weather conditions cause small variations in atmospheric pressure and therefore in boiling temperature. The corresponding change in boiling temperature is generally no more than about $1^\circ\mathrm{C}$.

\textbf{\\ Some Consequences of $T_{\mathrm{sat}}$ and $P_{\mathrm{sat}}$ Dependence}

Since the boiling temperature is determined by pressure, boiling can be controlled by controlling pressure.

This dependence has several practical applications.

\textbf{\\ Refrigerant-134a}

Consider a sealed can of liquid refrigerant-134a at $25^\circ\mathrm{C}$.

When the lid is slowly opened, refrigerant escapes and the pressure inside the can decreases.

As the pressure falls to $1~\mathrm{atm}$, the saturation temperature of refrigerant-134a is approximately:

\begin{equation*}
T_{\mathrm{sat}}=-26^\circ\mathrm{C}
\end{equation*}

The refrigerant temperature therefore drops rapidly and may cause ice to form on the outside of the can when the surrounding air is humid.

The liquid refrigerant remains at approximately $-26^\circ\mathrm{C}$ until the last drop vaporizes.

At $1~\mathrm{atm}$, the latent heat of vaporization of refrigerant-134a is:

\begin{equation*}
h_{\mathrm{vap}}=217~\mathrm{kJ/kg}
\end{equation*}

A liquid cannot vaporize unless it absorbs the required latent heat of vaporization.

Therefore, the rate of vaporization depends on the rate of heat transfer to the refrigerant:

\begin{equation*}
\boxed{\text{Higher heat-transfer rate}
\quad\Rightarrow\quad
\text{Higher vaporization rate}}
\end{equation*}

Heavy insulation reduces heat transfer and therefore reduces the rate of vaporization.

In the limiting case of no heat transfer, the refrigerant can remain as a liquid at $-26^\circ\mathrm{C}$ indefinitely.

\textbf{\\ Liquid Nitrogen}

The boiling temperature of nitrogen at atmospheric pressure is approximately:

\begin{equation*}
T_{\mathrm{sat}}=-196^\circ\mathrm{C}
\end{equation*}

Liquid nitrogen exposed to the atmosphere therefore remains at approximately $-196^\circ\mathrm{C}$ while it is evaporating.

Its temperature remains nearly constant until the liquid nitrogen is depleted.

This property makes liquid nitrogen useful for:
\begin{itemize}
    \item low-temperature scientific studies,
    \item superconductivity experiments, and
    \item cryogenic applications.
\end{itemize}

\textbf{\\ Vacuum Cooling}

\textbf{Vacuum cooling} is a method of cooling leafy vegetables by reducing the pressure of a sealed cooling chamber.

The pressure is reduced to the saturation pressure corresponding to the desired low temperature. Some water then evaporates from the products.

The latent heat required for vaporization is absorbed from the products, lowering their temperature.

For water at $0^\circ\mathrm{C}$:

\begin{equation*}
P_{\mathrm{sat}}=0.61~\mathrm{kPa}
\end{equation*}

Therefore, products can be cooled to approximately $0^\circ\mathrm{C}$ by reducing the chamber pressure to about $0.61~\mathrm{kPa}$.

Reducing the pressure below $0.61~\mathrm{kPa}$ can increase the cooling rate, but this is undesirable because of:
\begin{itemize}
    \item the possibility of freezing, and
    \item the additional cost.
\end{itemize}

\textbf{\\ Vacuum Freezing}

Vacuum cooling becomes \textbf{vacuum freezing} when the vapor pressure in the vacuum chamber is reduced below the saturation pressure of water at $0^\circ\mathrm{C}$.

\begin{equation*}
P<0.61~\mathrm{kPa}
\quad\Rightarrow\quad
\text{Vacuum freezing}
\end{equation*}

The reduced pressure causes water to evaporate and can eventually produce ice.

The principle of making ice using a vacuum pump was demonstrated by Dr.~William Cullen in Scotland in 1775 by evacuating the airspace in a water tank.

\textbf{\\ Package Icing}

\textbf{Package icing} is used in small-scale cooling applications to remove heat and keep products cool during transportation.

It takes advantage of the large latent heat of fusion of water.

Package icing is suitable only for products that are not harmed by direct contact with ice.

In addition to refrigeration, ice also provides moisture.

\textbf{\\ Summary of Phase-Change States}

\begin{center}
    \begin{tabular}{|c|c|c|}
        \hline
        \textbf{State} & \textbf{Definition} & \textbf{Characteristic} \\
        \hline
        Compressed liquid & Liquid not about to vaporize & $T<T_{\mathrm{sat}}$ at specified $P$ \\
        \hline
        Saturated liquid & Liquid about to vaporize & Beginning of vaporization \\
        \hline
        Saturated mixture & Liquid and vapor coexist & Two-phase equilibrium \\
        \hline
        Saturated vapor & Vapor about to condense & End of vaporization \\
        \hline
        Superheated vapor & Vapor not about to condense & $T>T_{\mathrm{sat}}$ at specified $P$ \\
        \hline
    \end{tabular}
\end{center}





























\mysection{3.4}{Property Diagrams for Phase-Change Processes}{3.4}

The variations of properties during phase-change processes are best studied and understood with the help of \textbf{property diagrams}.

For pure substances, the important property diagrams are:
\begin{itemize}
    \item $T$-$v$ diagram,
    \item $P$-$v$ diagram, and
    \item $P$-$T$ diagram.
\end{itemize}

The three-dimensional representation of the thermodynamic behavior of a substance is called the \textbf{$P$-$v$-$T$ surface}.

\textbf{\\ 1. The $T$-$v$ Diagram}

The $T$-$v$ diagram represents the relationship between temperature $T$ and specific volume $v$ during phase-change processes.

Consider the heating of water at constant pressure.

As pressure is increased:
\begin{itemize}
    \item the saturation temperature increases,
    \item the specific volume of the saturated liquid increases,
    \item the specific volume of the saturated vapor decreases, and
    \item the horizontal distance between the saturated liquid and saturated vapor states decreases.
\end{itemize}

The saturation line therefore becomes progressively shorter as pressure increases.

\textbf{\\ Critical Point}

As the pressure is increased further, the saturated-liquid and saturated-vapor states approach each other.

At a certain pressure, the two states become identical. This point is called the \textbf{critical point}.

The critical point is defined as the point at which the saturated liquid and saturated vapor states are identical.

The properties at the critical point are called:
\begin{itemize}
    \item critical temperature, $T_{\mathrm{cr}}$,
    \item critical pressure, $P_{\mathrm{cr}}$, and
    \item critical specific volume, $v_{\mathrm{cr}}$.
\end{itemize}

\insertfigure{0.6}{Figure : 5.png}

\textbf{\\ Saturated Liquid and Saturated Vapor Lines}

The saturated liquid states at different pressures can be connected to form the \textbf{saturated liquid line}.

Similarly, the saturated vapor states can be connected to form the \textbf{saturated vapor line}.

These two lines meet at the critical point and form a dome-shaped region.

\begin{equation*}
\boxed{\text{Saturated liquid line}+\text{Saturated vapor line}
\rightarrow \text{Saturation dome}}
\end{equation*}

\insertfigure{0.6}{Figure : 6.png}

\textbf{\\ Regions of the $T$-$v$ Diagram}

The $T$-$v$ diagram is divided into three principal regions:

\begin{center}
    \begin{tabular}{|c|c|c|}
        \hline
        \textbf{Region} & \textbf{Location} & \textbf{Phase} \\
        \hline
        Compressed liquid region & Left of saturated liquid line & Liquid \\
        \hline
        Saturated liquid--vapor region & Under the dome & Liquid + Vapor \\
        \hline
        Superheated vapor region & Right of saturated vapor line & Vapor \\
        \hline
    \end{tabular}
\end{center}

The compressed liquid and superheated vapor regions are \textbf{single-phase regions}.

The region under the dome is a \textbf{two-phase equilibrium region}, also called the \textbf{saturated liquid--vapor mixture region} or \textbf{wet region}.

\textbf{\\ Supercritical Region}

At pressures greater than the critical pressure:

\begin{equation*}
P>P_{\mathrm{cr}}
\end{equation*}

there is no distinct phase-change or boiling process.

The specific volume increases continuously as heat is added, and only one phase is present at every state.

Eventually, the substance resembles a vapor, but there is no distinct point at which the liquid changes into vapor.

Above the critical state, there is no line separating the compressed liquid region from the superheated vapor region.

By convention:
\begin{itemize}
    \item at temperatures above $T_{\mathrm{cr}}$, the substance is referred to as superheated vapor,
    \item at temperatures below $T_{\mathrm{cr}}$, it is referred to as compressed liquid.
\end{itemize}

\textbf{\\ 2. The $P$-$v$ Diagram}

The general shape of the $P$-$v$ diagram is similar to that of the $T$-$v$ diagram.

The major difference is that the constant-temperature lines on the $P$-$v$ diagram have a \textbf{downward trend}.

As pressure decreases:
\begin{itemize}
    \item the specific volume increases slightly,
    \item the pressure eventually reaches the saturation pressure corresponding to $150^\circ\mathrm{C}$,
    \item boiling begins, and
    \item \textbf{during vaporization, pressure and temperature remain constant while specific volume increases.}
\end{itemize}

Once the last drop of liquid has vaporized, further reduction in pressure causes the specific volume to increase further.

\insertfigure{0.6}{Figure : 7.png}

\textbf{\\ Regions of the $P$-$v$ Diagram}

The $P$-$v$ diagram contains the same basic regions as the $T$-$v$ diagram:

\begin{center}
    \begin{tabular}{|c|c|}
        \hline
        \textbf{Region} & \textbf{Phase} \\
        \hline
        Compressed liquid region & Liquid \\
        \hline
        Saturated liquid--vapor region & Liquid + Vapor \\
        \hline
        Superheated vapor region & Vapor \\
        \hline
    \end{tabular}
\end{center}

The saturated liquid line and saturated vapor line meet at the critical point.

\textbf{\\ 3. Extending the Diagrams to Include the Solid Phase}

The $T$-$v$ and $P$-$v$ diagrams initially describe equilibrium states involving the liquid and vapor phases.

They can also be extended to include:
\begin{itemize}
    \item the solid phase,
    \item the solid--liquid equilibrium region, and
    \item the solid--vapor equilibrium region.
\end{itemize}

The same basic principles that apply to liquid--vapor phase changes also apply to solid--liquid and solid--vapor phase changes.

\textbf{\\ Behavior During Freezing}

Most substances \textbf{contract during freezing}.

Water is an important exception because it \textbf{expands during freezing}.

The $P$-$v$ diagrams for these two classes of substances differ mainly in the solid--liquid saturation region.

The corresponding $T$-$v$ diagrams are similar to the $P$-$v$ diagrams, especially for substances that contract during freezing.


\insertfigure{1.0}{Figure : 8.png}

\textbf{\\ Importance of Water's Expansion During Freezing}

Water's expansion during freezing has important consequences in nature.

If water contracted during freezing, as most substances do, the resulting ice would be denser than liquid water.

The ice would then sink to the bottom of rivers, lakes, and oceans instead of floating on the surface.

The resulting ice layers would prevent sunlight from reaching the lower regions of the water bodies, potentially covering their bottoms with ice and seriously disrupting marine life.

\textbf{\\ Triple-Phase Equilibrium}

Under certain conditions, all three phases of a pure substance can coexist in equilibrium:

\begin{equation*}
\boxed{\text{Solid}+\text{Liquid}+\text{Vapor}}
\end{equation*}

On the $P$-$v$ and $T$-$v$ diagrams, these three-phase equilibrium states form a line called the \textbf{triple line}.

The states along the triple line have:
\begin{itemize}
    \item the same pressure,
    \item the same temperature, and
    \item different specific volumes.
\end{itemize}

On a $P$-$T$ diagram, the triple line appears as a point. Therefore, it is called the \textbf{triple point}.

For water:

\begin{equation*}
T_{\mathrm{tp}}=0.01^\circ\mathrm{C}
\end{equation*}

\begin{equation*}
P_{\mathrm{tp}}=0.6117~\mathrm{kPa}
\end{equation*}

Thus, all three phases of water coexist in equilibrium only at precisely:

\begin{equation*}
\boxed{
T=0.01^\circ\mathrm{C},
\qquad
P=0.6117~\mathrm{kPa}
}
\end{equation*}

\textbf{\\ Triple Line and Triple Point}

\begin{center}
    \begin{tabular}{|c|c|c|}
        \hline
        \textbf{Diagram} & \textbf{Representation} & \textbf{Name} \\
        \hline
        $P$-$v$ diagram & Line & Triple line \\
        \hline
        $T$-$v$ diagram & Line & Triple line \\
        \hline
        $P$-$T$ diagram & Point & Triple point \\
        \hline
    \end{tabular}
\end{center}

\textbf{\\ Sublimation}

A substance can pass directly from the solid phase to the vapor phase without becoming a liquid.

This process is called \textbf{sublimation}.

There are two possible paths from solid to vapor:

\begin{equation*}
\text{Solid}
\rightarrow
\text{Liquid}
\rightarrow
\text{Vapor}
\end{equation*}

or

\begin{equation*}
\boxed{\text{Solid}\rightarrow\text{Vapor}}
\end{equation*}

The direct solid-to-vapor transition occurs at pressures below the triple-point pressure because a pure substance cannot exist as a liquid in stable equilibrium at those pressures.

For substances whose triple-point pressure is greater than atmospheric pressure, sublimation is the only way to change directly from solid to vapor at atmospheric conditions.

A common example is solid carbon dioxide, $\mathrm{CO_2}$, known as \textbf{dry ice}.

\textbf{\\ 4. The $P$-$T$ Diagram}

The $P$-$T$ diagram is often called the \textbf{phase diagram} because it shows the regions in which the solid, liquid, and vapor phases exist.

The three phases are separated by three equilibrium lines:

\begin{center}
    \begin{tabular}{|c|c|c|}
        \hline
        \textbf{Line} & \textbf{Phases Separated} & \textbf{Phase Change} \\
        \hline
        Sublimation line & Solid and Vapor & Sublimation \\
        \hline
        Melting line & Solid and Liquid & Melting/Fusion \\
        \hline
        Vaporization line & Liquid and Vapor & Vaporization \\
        \hline
    \end{tabular}
\end{center}

All three lines meet at the triple point.

\begin{equation*}
\boxed{
\text{Sublimation line}
\cap
\text{Melting line}
\cap
\text{Vaporization line}
=
\text{Triple point}
}
\end{equation*}

The vaporization line terminates at the critical point.

Beyond the critical point, there is no distinction between the liquid and vapor phases.

Therefore:

\begin{equation*}
\boxed{\text{Vaporization line ends at the critical point}}
\end{equation*}

\textbf{\\ Substances That Contract and Expand on Freezing}

Substances that contract on freezing and substances that expand on freezing differ only in the behavior of their \textbf{melting line} on the $P$-$T$ diagram.

\begin{center}
    \begin{tabular}{|c|c|}
        \hline
        \textbf{Behavior} & \textbf{Melting Line} \\
        \hline
        Substance contracts on freezing & Positive slope \\
        \hline
        Substance expands on freezing & Negative slope \\
        \hline
    \end{tabular}
\end{center}

Water belongs to the second category because it expands when it freezes.

\insertfigure{0.6}{Figure : 9.png}

\textbf{\\ 5. The $P$-$v$-$T$ Surface}

The complete thermodynamic behavior of a pure substance can be represented using a three-dimensional \textbf{$P$-$v$-$T$ surface}.

The state of a simple compressible substance is fixed by any two independent intensive properties.

Once two appropriate independent properties are specified, all other properties become dependent properties.

If a property relation is written in the general form:

\begin{equation*}
z=f(x,y)
\end{equation*}

it represents a surface in three-dimensional space.

For the $P$-$v$-$T$ surface:
\begin{itemize}
    \item $T$ and $v$ can be considered independent variables,
    \item $P$ is the dependent variable.
\end{itemize}

Thus, $T$ and $v$ form the base of the surface while $P$ represents the height.

\textbf{\\ Equilibrium States on the $P$-$v$-$T$ Surface}

Every point on the $P$-$v$-$T$ surface represents an \textbf{equilibrium state}.

A quasi-equilibrium process passes through a sequence of equilibrium states. Therefore, every state along a quasi-equilibrium process lies on the $P$-$v$-$T$ surface.

\textbf{\\ Regions on the $P$-$v$-$T$ Surface}

The single-phase regions appear as curved surfaces on the $P$-$v$-$T$ surface.

The two-phase regions appear as surfaces perpendicular to the $P$-$T$ plane.

This occurs because the projections of the two-phase regions onto the $P$-$T$ plane are lines.

\textbf{\\ Relationship Between Property Diagrams}

The commonly used two-dimensional property diagrams are projections of the three-dimensional $P$-$v$-$T$ surface.

\begin{center}
    \begin{tabular}{|c|c|}
        \hline
        \textbf{Diagram} & \textbf{Relationship to $P$-$v$-$T$ Surface} \\
        \hline
        $P$-$v$ diagram & Projection onto the $P$-$v$ plane \\
        \hline
        $T$-$v$ diagram & Bird's-eye view of the surface \\
        \hline
        $P$-$T$ diagram & Projection onto the $P$-$T$ plane \\
        \hline
    \end{tabular}
\end{center}

\insertfigure{0.6}{Figure : 10.png}

\textbf{\\ Critical Point vs Triple Point}

\begin{center}
    \begin{tabular}{|c|c|c|}
        \hline
        \textbf{Property} & \textbf{Critical Point} & \textbf{Triple Point} \\
        \hline
        Phases involved & Liquid + Vapor become identical & Solid + Liquid + Vapor coexist \\
        \hline
        Location & End of vaporization line & Intersection of three phase lines \\
        \hline
        $P$-$T$ diagram & Point & Point \\
        \hline
        $T$-$v$ diagram & Point & Triple line \\
        \hline
        $P$-$v$ diagram & Point & Triple line \\
        \hline
    \end{tabular}
\end{center}



























\mysection{3.5}{Property Tables}{3.5}

For most substances, the relationships among thermodynamic properties are too complex to be expressed by simple equations. Therefore, thermodynamic properties are frequently presented in the form of \textbf{property tables}.

Some thermodynamic properties can be measured directly, while others are calculated using relationships between measurable properties. The results are compiled into tables in a convenient form.

For each substance, separate tables are generally provided for different regions, such as:
\begin{itemize}
    \item saturated liquid and saturated vapor,
    \item saturated liquid--vapor mixture,
    \item superheated vapor, and
    \item compressed liquid.
\end{itemize}

Property tables may be provided in different unit systems.

\textbf{\\ Enthalpy---A Combination Property}

In many thermodynamic processes, particularly in power generation and refrigeration, the combination

\begin{equation*}
u+Pv
\end{equation*}

occurs frequently.

For convenience, this combination is defined as a new thermodynamic property called \textbf{enthalpy}, denoted by $h$.

The specific enthalpy is defined as

\begin{equation*}
\boxed{h=u+Pv}
\end{equation*}

The total enthalpy is

\begin{equation*}
\boxed{H=U+PV}
\end{equation*}


The pressure--volume product has energy units. Therefore, the definition of enthalpy is dimensionally consistent.

If internal energy is not listed in a property table, it can be obtained from

\begin{equation*}
\boxed{u=h-Pv}
\end{equation*}

The widespread use of enthalpy is associated with its importance in the analysis of control-volume devices such as turbines, compressors, and refrigeration systems.

\textbf{\\ 1a. Saturated Liquid and Saturated Vapor States}

The properties of saturated liquid and saturated vapor are commonly provided in saturation tables.

Two forms of saturation tables are generally available:
\begin{itemize}
    \item tables arranged according to temperature, and
    \item tables arranged according to pressure.
\end{itemize}

Therefore:
\begin{itemize}
    \item use the temperature-based saturation table when temperature is given,
    \item use the pressure-based saturation table when pressure is given.
\end{itemize}

\textbf{\\ Subscript Notation}

The subscript $f$ denotes a \textbf{saturated liquid} property.

The subscript $g$ denotes a \textbf{saturated vapor} property.

Thus,

\begin{equation*}
v_f=\text{specific volume of saturated liquid}
\end{equation*}

\begin{equation*}
v_g=\text{specific volume of saturated vapor}
\end{equation*}

The subscript $fg$ denotes the difference between the saturated vapor and saturated liquid values of a property.

For specific volume,

\begin{equation*}
\boxed{v_{fg}=v_g-v_f}
\end{equation*}

Similarly,

\begin{equation*}
\boxed{u_{fg}=u_g-u_f}
\end{equation*}

and

\begin{equation*}
\boxed{h_{fg}=h_g-h_f}
\end{equation*}
    
\textbf{\\ Enthalpy of Vaporization}

The quantity $h_{fg}$ is called the \textbf{enthalpy of vaporization} or \textbf{latent heat of vaporization}.

It represents the energy required to vaporize a unit mass of saturated liquid at a specified temperature or pressure.

It is given by

\begin{equation*}
\boxed{h_{fg}=h_g-h_f}
\end{equation*}

The enthalpy of vaporization decreases as temperature or pressure increases and becomes zero at the critical point.

\insertfigure{1.0}{Figure : 11.png}

\textbf{\\ 1b. Saturated Liquid--Vapor Mixture}

During vaporization, a substance exists partly as saturated liquid and partly as saturated vapor.

Such a state is called a \textbf{saturated liquid--vapor mixture}.

To analyze the mixture, the relative amounts of the liquid and vapor phases must be determined.

This is done using a property called \textbf{quality}.

\textbf{\\ Quality}

Quality, denoted by $x$, is defined as the ratio of the mass of vapor to the total mass of the saturated mixture:

\begin{equation*}
\boxed{x=\frac{m_g}{m_t}}
\end{equation*}

where

\begin{equation*}
m_t=m_f+m_g
\end{equation*}

Here:
\begin{itemize}
    \item $m_f$ is the mass of saturated liquid,
    \item $m_g$ is the mass of saturated vapor, and
    \item $m_t$ is the total mass of the mixture.
\end{itemize}

Quality has meaning \textbf{only for saturated liquid--vapor mixtures}.

It has no meaning in:
\begin{itemize}
    \item the compressed liquid region, or
    \item the superheated vapor region.
\end{itemize}

For saturated mixtures,

\begin{equation*}
\boxed{0\leq x\leq 1}
\end{equation*}

The limiting cases are

\begin{equation*}
x=0
\quad\Rightarrow\quad
\text{saturated liquid}
\end{equation*}

\begin{equation*}
x=1
\quad\Rightarrow\quad
\text{saturated vapor}
\end{equation*}

Thus, quality specifies the fraction of the total mass that exists as vapor.

In a saturated mixture, quality can serve as one of the two independent intensive properties required to specify the state.

\textbf{\\ Average Properties of a Saturated Mixture}

A saturated liquid--vapor mixture can be treated as a combination of two subsystems:
\begin{itemize}
    \item saturated liquid, and
    \item saturated vapor.
\end{itemize}

The average specific volume of the mixture is

\begin{equation*}
v=(1-x)v_f+xv_g
\end{equation*}

Using

\begin{equation*}
v_{fg}=v_g-v_f
\end{equation*}

the equation becomes

\begin{equation*}
\boxed{v=v_f+xv_{fg}}
\end{equation*}

Therefore, quality can be determined from

\begin{equation*}
\boxed{x=\frac{v-v_f}{v_{fg}}}
\end{equation*}

Similarly, the specific internal energy of a saturated mixture is

\begin{equation*}
\boxed{u=u_f+xu_{fg}}
\end{equation*}

and the specific enthalpy is

\begin{equation*}
\boxed{h=h_f+xh_{fg}}
\end{equation*}

These relations can be expressed in a general form as

\begin{equation*}
\boxed{y=y_f+xy_{fg}}
\end{equation*}

where $y$ can represent any of the properties

\begin{equation*}
y=v,\quad u,\quad h
\end{equation*}

The average mixture property always lies between the corresponding saturated-liquid and saturated-vapor values:

\begin{equation*}
\boxed{y_f\leq y\leq y_g}
\end{equation*}

Hence,

\begin{equation*}
v_f\leq v\leq v_g
\end{equation*}

\begin{equation*}
u_f\leq u\leq u_g
\end{equation*}

\begin{equation*}
h_f\leq h\leq h_g
\end{equation*}

\textbf{\\ Quality and Property Diagrams}

On a $P$-$v$ or $T$-$v$ diagram, quality is related to the horizontal distance between the saturated liquid and saturated vapor states.

At a specified pressure or temperature,

\begin{equation*}
\boxed{
x=
\frac{\text{distance from saturated liquid state to actual state}}
{\text{distance from saturated liquid state to saturated vapor state}}
}
\end{equation*}

Therefore:
\begin{itemize}
    \item a state close to the saturated liquid line has low quality,
    \item a state close to the saturated vapor line has high quality,
    \item quality increases from $0$ at saturated liquid to $1$ at saturated vapor.
\end{itemize}

All saturated-mixture states lie under the saturation dome.

To determine the properties of a saturated mixture, saturated-liquid and saturated-vapor data are required.

\textbf{\\ Saturated Solid--Vapor Mixtures}

Property tables are also available for saturated solid--vapor mixtures.

Saturated solid--vapor mixtures can be analyzed in the same general manner as saturated liquid--vapor mixtures.

\textbf{\\ 2. Superheated Vapor}

A substance exists as \textbf{superheated vapor} when it lies in the superheated vapor region.

The superheated region is a \textbf{single-phase region} containing only vapor.

Therefore, temperature and pressure can be used as the two independent properties for superheated vapor.

Superheated vapor tables generally list properties as functions of temperature for selected pressures.

The saturation temperature corresponding to each pressure is also indicated in the tables.

\textbf{\\ Characteristics of Superheated Vapor}

Compared with saturated vapor, superheated vapor is characterized by:

\begin{equation*}
\boxed{P<P_{\mathrm{sat}}\quad\text{at a given }T}
\end{equation*}

\begin{equation*}
\boxed{T>T_{\mathrm{sat}}\quad\text{at a given }P}
\end{equation*}

\begin{equation*}
\boxed{v>v_g\quad\text{at a given }P\text{ or }T}
\end{equation*}

\begin{equation*}
\boxed{u>u_g\quad\text{at a given }P\text{ or }T}
\end{equation*}

\begin{equation*}
\boxed{h>h_g\quad\text{at a given }P\text{ or }T}
\end{equation*}

Thus, superheated vapor generally has higher specific volume, internal energy, and enthalpy than saturated vapor at the same pressure or temperature.

\textbf{\\ Identifying Superheated Vapor}

At a specified pressure:

\begin{equation*}
T>T_{\mathrm{sat}}
\quad\Rightarrow\quad
\text{superheated vapor}
\end{equation*}

At a specified temperature:

\begin{equation*}
P<P_{\mathrm{sat}}
\quad\Rightarrow\quad
\text{superheated vapor}
\end{equation*}

Quality is not defined in the superheated vapor region.

\textbf{\\ 3. Compressed Liquid}

Compressed liquid tables are less commonly available than superheated vapor tables.

The properties of compressed liquids vary only mildly with pressure.

In general, the properties of a compressed liquid depend much more strongly on temperature than on pressure.

Therefore, when compressed-liquid data are unavailable, a compressed liquid can often be approximated as a saturated liquid at the same temperature.

For a compressed liquid,

\begin{equation*}
\boxed{y\approx y_f@T}
\end{equation*}

where $y$ may represent

\begin{equation*}
y=v,\quad u,\quad h
\end{equation*}

Thus,

\begin{equation*}
\boxed{v\approx v_f@T}
\end{equation*}

\begin{equation*}
\boxed{u\approx u_f@T}
\end{equation*}

\begin{equation*}
\boxed{h\approx h_f@T}
\end{equation*}

This approximation is generally quite good for $v$ and $u$.

The enthalpy $h$ is more sensitive to pressure, so the approximation may introduce a more significant error for enthalpy.

\textbf{\\ Improved Approximation for Enthalpy}

At low to moderate pressures and temperatures, the enthalpy of a compressed liquid can be estimated using

\begin{equation*}
\boxed{
h\approx h_f@T+v_f@T
\left(P-P_{\mathrm{sat}}@T\right)
}
\end{equation*}

This approximation may significantly improve the estimate of enthalpy under appropriate conditions.

At moderate to high temperatures and pressures, however, the approximation may not provide significant improvement and can even introduce greater error.

\textbf{\\ Characteristics of Compressed Liquid}

Compared with saturated liquid, a compressed liquid is characterized by:

\begin{equation*}
\boxed{P>P_{\mathrm{sat}}\quad\text{at a given }T}
\end{equation*}

\begin{equation*}
\boxed{T<T_{\mathrm{sat}}\quad\text{at a given }P}
\end{equation*}

\begin{equation*}
\boxed{v<v_f\quad\text{at a given }P\text{ or }T}
\end{equation*}

\begin{equation*}
\boxed{u<u_f\quad\text{at a given }P\text{ or }T}
\end{equation*}

\begin{equation*}
\boxed{h<h_f\quad\text{at a given }P\text{ or }T}
\end{equation*}

Unlike superheated vapor, the properties of a compressed liquid are generally not very different from the corresponding saturated-liquid values.

Quality has no meaning in the compressed liquid region.

\textbf{\\ Identifying the Thermodynamic Region}

When the state of a pure substance is not immediately known, the appropriate saturation data can be used to determine the region.

If pressure and temperature are given, compare the specified temperature with the saturation temperature at the specified pressure:

\begin{equation*}
\boxed{
T<T_{\mathrm{sat}}@P
\quad\Rightarrow\quad
\text{compressed liquid}
}
\end{equation*}

\begin{equation*}
\boxed{
T=T_{\mathrm{sat}}@P
\quad\Rightarrow\quad
\text{saturated state}
}
\end{equation*}

\begin{equation*}
\boxed{
T>T_{\mathrm{sat}}@P
\quad\Rightarrow\quad
\text{superheated vapor}
}
\end{equation*}

Alternatively, when temperature is given, compare pressure with the saturation pressure at that temperature:

\begin{equation*}
\boxed{
P>P_{\mathrm{sat}}@T
\quad\Rightarrow\quad
\text{compressed liquid}
}
\end{equation*}

\begin{equation*}
\boxed{
P=P_{\mathrm{sat}}@T
\quad\Rightarrow\quad
\text{saturated state}
}
\end{equation*}

\begin{equation*}
\boxed{
P<P_{\mathrm{sat}}@T
\quad\Rightarrow\quad
\text{superheated vapor}
}
\end{equation*}

The region can also be determined when internal energy, enthalpy, or specific volume is known by comparing the specified property with its saturated-liquid and saturated-vapor values.

For a generic property $y$:

\begin{equation*}
\boxed{
y<y_f
\quad\Rightarrow\quad
\text{compressed liquid}
}
\end{equation*}

\begin{equation*}
\boxed{
y_f\leq y\leq y_g
\quad\Rightarrow\quad
\text{saturated liquid--vapor mixture}
}
\end{equation*}

\begin{equation*}
\boxed{
y>y_g
\quad\Rightarrow\quad
\text{superheated vapor}
}
\end{equation*}

This criterion can be applied to properties such as $u$, $h$, and $v$.

\textbf{\\ Reference State and Reference Values}

The absolute values of some thermodynamic properties cannot be measured directly.

Properties such as internal energy, enthalpy, and entropy are calculated from relationships involving measurable properties.

These relationships determine \textbf{changes in properties}, rather than absolute values.

Therefore, a convenient \textbf{reference state} is selected and one or more properties are assigned zero values at that state.

For water, the reference state is chosen as a particular saturated-liquid state, with internal energy and entropy assigned zero values there.

For refrigerant-134a, a different saturated-liquid reference state is selected, with enthalpy and entropy assigned zero values there.

The choice of reference state can cause different property tables to list different absolute values for some properties.

This does not cause a problem in thermodynamic calculations because calculations generally involve \textbf{changes in properties}.

As long as values are taken from one consistent set of tables or charts, the selected reference state has no effect on the calculated property changes.

Some properties may have negative values because of the selected reference state.

\textbf{\\ Consistency of Property Tables}

Different property tables may use different reference states.

Therefore:

\begin{equation*}
\boxed{
\text{Always use values from one consistent set of tables or charts.}
}
\end{equation*}

Absolute values may differ between table sets, but consistent property differences remain suitable for thermodynamic calculations.

\textbf{\\ Quick Property-Table Selection Guide}

\begin{center}
\begin{tabular}{|c|c|}
\hline
\textbf{Known Information / Region} & \textbf{Appropriate Table} \\
\hline
Saturated state, temperature given & Temperature-based saturation table \\
\hline
Saturated state, pressure given & Pressure-based saturation table \\
\hline
Superheated vapor & Superheated vapor table \\
\hline
Compressed liquid & Compressed liquid table \\
\hline
Saturated solid--vapor mixture & Saturated solid--vapor table \\
\hline
\end{tabular}
\end{center}



















\mysection{3.6}{The Ideal-Gas Equation of State}{3.6}

Property tables provide accurate information about thermodynamic properties, but they can be bulky and may contain typographical errors. A more practical approach is to use simple relationships among thermodynamic properties whenever they provide sufficient accuracy.

\textbf{\\ Equation of State}

An \textbf{equation of state} is an equation that relates the pressure, temperature, and specific volume of a substance.

More generally, property relationships involving other properties of a substance at equilibrium states may also be referred to as equations of state.

Several equations of state exist, ranging from simple to very complex.

The simplest and best-known equation of state for substances in the gas phase is the \textbf{ideal-gas equation of state}.

It predicts the $P$-$v$-$T$ behavior of gases accurately within an appropriately selected range.

\textbf{\\ Gas and Vapor}

The terms \textbf{gas} and \textbf{vapor} are often used synonymously.

However, there is a conventional distinction:
\begin{itemize}
        \item The term \textbf{vapor} usually implies a gas that is not far from a state of condensation.
\end{itemize}

\textbf{\\ Experimental Basis of the Ideal-Gas Relation}

The ideal-gas relation is based on experimentally observed behavior.

\textbf{Boyle's Law}

Robert Boyle observed that, at constant temperature, the pressure of a gas is inversely proportional to its volume.

Thus,

\begin{equation*}
P\propto\frac{1}{V}
\end{equation*}

for a fixed amount of gas at constant temperature.

\textbf{\\ Charles's and Gay-Lussac's Observations}

J.~Charles and J.~Gay-Lussac experimentally determined that, at low pressures, the volume of a gas is proportional to its absolute temperature.

Thus,

\begin{equation*}
V\propto T
\end{equation*}

for a fixed amount of gas at constant pressure.

Combining these observations gives the ideal-gas relationship.

\textbf{\\ Ideal-Gas Equation of State}

The ideal-gas equation of state is

\begin{equation*}
\boxed{Pv=RT}
\end{equation*}

where:
\begin{itemize}
    \item $P$ is the \textbf{absolute pressure},
    \item $T$ is the \textbf{absolute temperature},
    \item $v$ is the \textbf{specific volume}, and
    \item $R$ is the \textbf{gas constant}.
\end{itemize}

A gas that obeys this relation is called an \textbf{ideal gas}.

The temperature in the ideal-gas equation must be an \textbf{absolute temperature}.

\textbf{\\ Gas Constant}

The gas constant $R$ is different for each gas.

It is determined from

\begin{equation*}
\boxed{R=\frac{R_u}{M}}
\end{equation*}

where:
\begin{itemize}
    \item $R_u$ is the \textbf{universal gas constant}, and
    \item $M$ is the \textbf{molar mass} of the gas.
\end{itemize}

Therefore, different gases have different specific gas constants even though they have the same universal gas constant.

\textbf{\\ Universal Gas Constant}

The universal gas constant $R_u$ has the same value for all substances.

Its numerical value depends only on the unit system being used.

\textbf{\\ Molar Mass}

The \textbf{molar mass}, $M$, is the mass associated with a specified amount of a substance on a mole basis.

It may be expressed on a molar basis appropriate to the unit system being used.

The mass of a system is related to molar mass and mole number by

\begin{equation*}
\boxed{m=MN}
\end{equation*}

where:
\begin{itemize}
    \item $m$ is the total mass,
    \item $M$ is the molar mass, and
    \item $N$ is the number of moles.
\end{itemize}

The values of $R$ and $M$ for different substances are tabulated in thermodynamic property tables.

\textbf{\\ Alternative Forms of the Ideal-Gas Equation}

Starting from

\begin{equation*}
Pv=RT
\end{equation*}

and using

\begin{equation*}
V=mv
\end{equation*}

the ideal-gas relation can be written in terms of total quantities as

\begin{equation*}
\boxed{PV=mRT}
\end{equation*}

Using

\begin{equation*}
m=MN
\end{equation*}

and

\begin{equation*}
R=\frac{R_u}{M}
\end{equation*}

gives

\begin{equation*}
mR=(MN)\left(\frac{R_u}{M}\right)=NR_u
\end{equation*}

Therefore,

\begin{equation*}
\boxed{PV=NR_uT}
\end{equation*}

This form is particularly convenient when the amount of gas is specified in terms of the number of moles.

\textbf{\\ Molar Specific Volume}

The \textbf{molar specific volume}, denoted by $\bar{v}$, is the volume per unit mole.

It is related to the total volume and mole number by

\begin{equation*}
\bar{v}=\frac{V}{N}
\end{equation*}

The ideal-gas equation can therefore be written as

\begin{equation*}
\boxed{P\bar{v}=R_uT}
\end{equation*}

Properties expressed on a unit-mole basis are commonly distinguished by a bar over the property.

For example:

\begin{equation*}
\bar{v}=\text{molar specific volume}
\end{equation*}

while

\begin{equation*}
v=\text{specific volume per unit mass}
\end{equation*}

Similarly, molar internal energy and molar enthalpy may be distinguished from their mass-based counterparts using a bar.

\textbf{\\ Ideal-Gas Relation Between Two States}

For a fixed mass of an ideal gas, the gas constant remains constant.

Writing the ideal-gas equation for two states gives

\begin{equation*}
P_1V_1=mRT_1
\end{equation*}

and

\begin{equation*}
P_2V_2=mRT_2
\end{equation*}

Since $m$ and $R$ are constant,

\begin{equation*}
\boxed{
\frac{P_1V_1}{T_1}
=
\frac{P_2V_2}{T_2}
}
\end{equation*}

This relation is useful for analyzing processes involving a fixed mass of an ideal gas.

\textbf{\\ Special Cases for a Fixed Mass of Ideal Gas}

The two-state ideal-gas relation can be simplified when one of the properties remains constant.

\textbf{\\ Constant Volume}

If

\begin{equation*}
V=\text{constant}
\end{equation*}

then

\begin{equation*}
\boxed{\frac{P}{T}=\text{constant}}
\end{equation*}

Therefore,

\begin{equation*}
\boxed{\frac{P_1}{T_1}=\frac{P_2}{T_2}}
\end{equation*}

Thus, for a fixed amount of ideal gas at constant volume, pressure is directly proportional to absolute temperature.

\textbf{\\ Constant Pressure}

If

\begin{equation*}
P=\text{constant}
\end{equation*}

then

\begin{equation*}
\boxed{\frac{V}{T}=\text{constant}}
\end{equation*}

Therefore,

\begin{equation*}
\boxed{\frac{V_1}{T_1}=\frac{V_2}{T_2}}
\end{equation*}

Thus, for a fixed amount of ideal gas at constant pressure, volume is directly proportional to absolute temperature.

\textbf{\\ Constant Temperature}

If

\begin{equation*}
T=\text{constant}
\end{equation*}

then

\begin{equation*}
\boxed{PV=\text{constant}}
\end{equation*}

Therefore,

\begin{equation*}
\boxed{P_1V_1=P_2V_2}
\end{equation*}

Thus, for a fixed amount of ideal gas at constant temperature, pressure is inversely proportional to volume.

\textbf{\\ Characteristics of an Ideal Gas}

An ideal gas is an \textbf{imaginary substance} that obeys

\begin{equation*}
Pv=RT
\end{equation*}

Real gases can behave very closely to an ideal gas under appropriate conditions.

Experimentally, the ideal-gas relation provides a good approximation for real gases when the gas has a sufficiently \textbf{low density}.

Low density is generally associated with:
\begin{itemize}
    \item low pressure, and
    \item high temperature.
\end{itemize}

Under these conditions, intermolecular effects become relatively insignificant and the ideal-gas relation becomes a good approximation.

\textbf{\\ When Can a Real Gas Be Treated as an Ideal Gas?}

The ideal-gas model is most appropriate when the gas is at:
\begin{itemize}
    \item relatively low pressure, and
    \item relatively high temperature.
\end{itemize}

Under such conditions, many common gases can be treated as ideal gases with negligible error.

Examples of gases that can often be treated as ideal gases in practical applications include:
\begin{itemize}
    \item air,
    \item nitrogen,
    \item oxygen,
    \item hydrogen,
    \item helium,
    \item argon,
    \item neon, and
    \item carbon dioxide.
\end{itemize}

Even some heavier gases can be treated as ideal gases under appropriate conditions.

\textbf{\\ When Should the Ideal-Gas Model Not Be Used?}

Dense gases should not automatically be treated as ideal gases.

For example:
\begin{itemize}
    \item water vapor in steam power plants, and
    \item refrigerant vapor in refrigeration systems
\end{itemize}

may deviate significantly from ideal-gas behavior under conditions of practical interest.

For such substances, \textbf{property tables} should be used when the ideal-gas approximation is not sufficiently accurate.

\textbf{\\ Absolute Pressure Requirement}

The ideal-gas equation requires \textbf{absolute pressure}.

If a pressure is given as gauge pressure, it must first be converted to absolute pressure:

\begin{equation*}
\boxed{
P_{\mathrm{abs}}
=
P_{\mathrm{gage}}+P_{\mathrm{atm}}
}
\end{equation*}

Using gauge pressure directly in the ideal-gas equation leads to an incorrect result.

\textbf{\\ Absolute Temperature Requirement}

The ideal-gas equation requires \textbf{absolute temperature}.

When temperatures are provided on the Celsius scale, they must be converted to the absolute temperature scale before substitution into the ideal-gas equation.

The general conversion is

\begin{equation*}
\boxed{T(\mathrm{K})=T(^\circ\mathrm{C})+273.15}
\end{equation*}

For English units, the corresponding absolute temperature scale must be used.

\textbf{\\ Ideal-Gas Equation and Density}

Since

\begin{equation*}
v=\frac{1}{\rho}
\end{equation*}

the ideal-gas equation

\begin{equation*}
Pv=RT
\end{equation*}

can also be written as

\begin{equation*}
\boxed{P=\rho RT}
\end{equation*}

or

\begin{equation*}
\boxed{\rho=\frac{P}{RT}}
\end{equation*}

Thus, the density of an ideal gas increases with pressure and decreases with absolute temperature when the gas constant is fixed.

\textbf{\\ Specific Volume Form}

The ideal-gas equation can also be expressed directly in terms of density:

\begin{equation*}
\boxed{v=\frac{RT}{P}}
\end{equation*}

This shows that the specific volume:
\begin{itemize}
    \item increases with absolute temperature, and
    \item decreases with absolute pressure.
\end{itemize}





















\mysection{3.7}{Compressibility Factor---A Measure of Deviation from Ideal-Gas Behavior}{3.7}

The ideal-gas equation of state is simple and convenient to use. However, real gases can deviate significantly from ideal-gas behavior, particularly near the \textbf{saturation region} and the \textbf{critical point}.

The deviation of a real gas from ideal-gas behavior can be accounted for by introducing a correction factor called the \textbf{compressibility factor}, denoted by $Z$.

\textbf{\\ Compressibility Factor}

The compressibility factor is defined as

\begin{equation*}
\boxed{Z=\frac{Pv}{RT}}
\end{equation*}

Therefore, the real-gas equation of state can be written as

\begin{equation*}
\boxed{Pv=ZRT}
\end{equation*}

For an ideal gas,

\begin{equation*}
Pv=RT
\end{equation*}

Comparing the two relations gives

\begin{equation*}
\boxed{Z=1}
\end{equation*}

Thus, the compressibility factor represents the extent to which the behavior of a real gas differs from that predicted by the ideal-gas equation.

\textbf{\\ Physical Interpretation of $Z$}

The compressibility factor can also be expressed as

\begin{equation*}
\boxed{
Z=\frac{v_{\mathrm{actual}}}{v_{\mathrm{ideal}}}
}
\end{equation*}

where the ideal-gas specific volume at the same pressure and temperature is

\begin{equation*}
\boxed{
v_{\mathrm{ideal}}=\frac{RT}{P}
}
\end{equation*}

Therefore,

\begin{equation*}
Z=
\frac{v_{\mathrm{actual}}}
{RT/P}
\end{equation*}

The value of $Z$ indicates the deviation from ideal-gas behavior:

\begin{equation*}
\boxed{Z=1\quad\Rightarrow\quad\text{ideal-gas behavior}}
\end{equation*}

\begin{equation*}
\boxed{Z>1\quad\Rightarrow\quad
v_{\mathrm{actual}}>v_{\mathrm{ideal}}}
\end{equation*}

\begin{equation*}
\boxed{Z<1\quad\Rightarrow\quad
v_{\mathrm{actual}}<v_{\mathrm{ideal}}}
\end{equation*}

The farther $Z$ is from unity, the greater the deviation from ideal-gas behavior.

\textbf{\\ Real-Gas Behavior}

Real gases do not behave identically at the same temperature and pressure.

However, gases behave in a remarkably similar manner when their temperature and pressure are normalized with respect to their respective critical properties.

This leads to the \textbf{principle of corresponding states}.

\textbf{\\ Reduced Properties}

The pressure and temperature of a gas can be normalized using the critical pressure and critical temperature.

The \textbf{reduced pressure} is defined as

\begin{equation*}
\boxed{
P_R=\frac{P}{P_{\mathrm{cr}}}
}
\end{equation*}

The \textbf{reduced temperature} is defined as

\begin{equation*}
\boxed{
T_R=\frac{T}{T_{\mathrm{cr}}}
}
\end{equation*}

where:
\begin{itemize}
    \item $P_R$ is the reduced pressure,
    \item $T_R$ is the reduced temperature,
    \item $P_{\mathrm{cr}}$ is the critical pressure, and
    \item $T_{\mathrm{cr}}$ is the critical temperature.
\end{itemize}

Reduced properties are dimensionless.

\textbf{\\ Principle of Corresponding States}

The \textbf{principle of corresponding states} states that the compressibility factor of different gases is approximately the same at the same reduced pressure and reduced temperature.

Thus,

\begin{equation*}
\boxed{
Z=f(P_R,T_R)
}
\end{equation*}

This allows compressibility data for many different gases to be represented using a generalized chart.

The generalized compressibility chart is constructed by combining experimentally determined compressibility-factor data for various gases.

\textbf{\\ Generalized Compressibility Chart}

\insertfigure{1.0}{Figure : 12.png}

The generalized compressibility chart provides the compressibility factor $Z$ as a function of:

\begin{equation*}
P_R
\qquad\text{and}\qquad
T_R
\end{equation*}

Once $P_R$ and $T_R$ are known, the value of $Z$ can be obtained from the chart.

The actual specific volume can then be calculated from

\begin{equation*}
\boxed{
v=Z\frac{RT}{P}
}
\end{equation*}

or

\begin{equation*}
\boxed{
v=Zv_{\mathrm{ideal}}
}
\end{equation*}

Thus, the compressibility factor acts as a correction to the ideal-gas prediction.

\textbf{\\ Behavior at Very Low Pressure}

At very low reduced pressures,

\begin{equation*}
\boxed{P_R\ll1}
\end{equation*}

gases behave approximately as ideal gases regardless of temperature.

Therefore,

\begin{equation*}
\boxed{Z\rightarrow1}
\end{equation*}

as pressure approaches zero.

Hence,

\begin{equation*}
P\rightarrow0
\quad\Rightarrow\quad
\text{real-gas behavior}\rightarrow\text{ideal-gas behavior}
\end{equation*}

This is an important limiting behavior of real gases.

\textbf{\\ Behavior at High Reduced Temperature}

At sufficiently high reduced temperatures,

\begin{equation*}
\boxed{T_R>2}
\end{equation*}

ideal-gas behavior can generally be assumed with good accuracy over a wide range of pressures, except at very high reduced pressures.

Thus, increasing temperature relative to the critical temperature generally reduces the deviation from ideal-gas behavior.

\textbf{\\ Behavior Near the Critical Point}

The deviation of a gas from ideal-gas behavior is greatest in the vicinity of the critical point.

Therefore, the ideal-gas equation should be used with particular caution when a gas is near its critical state.

In the neighborhood of the critical point:

\begin{equation*}
\boxed{|Z-1|\ \text{can become large}}
\end{equation*}

Consequently, the generalized compressibility chart is particularly useful when significant nonideal behavior is expected.

\textbf{\\ Interpretation of the Generalized Compressibility Chart}

The generalized compressibility chart plots

\begin{equation*}
Z=\frac{Pv}{RT}
\end{equation*}

against reduced pressure $P_R$ for different values of reduced temperature $T_R$.

The chart demonstrates that different gases follow approximately the same trends when their states are expressed in terms of reduced properties.

The chart therefore provides a practical method for estimating real-gas behavior without requiring individual equations of state for every gas.

\textbf{\\ Pseudo-Reduced Specific Volume}

When pressure and specific volume, or temperature and specific volume, are known instead of pressure and temperature, direct use of the generalized compressibility chart may require trial-and-error calculations.

To avoid this difficulty, another reduced property called the \textbf{pseudo-reduced specific volume} is defined.

It is given by

\begin{equation*}
\boxed{
v_R=
\frac{v_{\mathrm{actual}}}
{RT_{\mathrm{cr}}/P_{\mathrm{cr}}}
}
\end{equation*}

The pseudo-reduced specific volume differs from $P_R$ and $T_R$ because it is defined using $T_{\mathrm{cr}}$ and $P_{\mathrm{cr}}$ rather than the critical specific volume.

Thus,

\begin{equation*}
\boxed{
v_R=
\frac{vP_{\mathrm{cr}}}
{RT_{\mathrm{cr}}}
}
\end{equation*}

where:
\begin{itemize}
    \item $v_R$ is the pseudo-reduced specific volume,
    \item $v$ is the actual specific volume,
    \item $P_{\mathrm{cr}}$ is the critical pressure,
    \item $T_{\mathrm{cr}}$ is the critical temperature, and
    \item $R$ is the gas constant.
\end{itemize}

\textbf{\\ Use of $v_R$ in the Compressibility Chart}

Lines of constant pseudo-reduced specific volume are included on the generalized compressibility chart.

Therefore, if $P$ and $v$ are known:
\begin{enumerate}
    \item calculate the reduced pressure $P_R$,
    \item calculate the pseudo-reduced specific volume $v_R$,
    \item locate the corresponding values on the generalized compressibility chart, and
    \item determine the required property from the chart.
\end{enumerate}

Similarly, if $T$ and $v$ are known, the chart can be used to determine the required pressure.

The inclusion of constant-$v_R$ lines eliminates the need for time-consuming trial-and-error iterations.

\textbf{\\ Determining Specific Volume Using $Z$}

When pressure and temperature are known, the procedure is:

\begin{enumerate}
    \item Determine the critical pressure $P_{\mathrm{cr}}$ and critical temperature $T_{\mathrm{cr}}$ of the gas.
    \item Calculate the reduced pressure:
    \begin{equation*}
    P_R=\frac{P}{P_{\mathrm{cr}}}
    \end{equation*}
    \item Calculate the reduced temperature:
    \begin{equation*}
    T_R=\frac{T}{T_{\mathrm{cr}}}
    \end{equation*}
    \item Use $P_R$ and $T_R$ to obtain $Z$ from the generalized compressibility chart.
    \item Calculate the actual specific volume from
    \begin{equation*}
    \boxed{
    v=Z\frac{RT}{P}
    }
    \end{equation*}
\end{enumerate}

\textbf{\\ Determining Pressure Using the Generalized Chart}

If temperature and specific volume are known, the procedure can be based on the pseudo-reduced specific volume.

First calculate

\begin{equation*}
\boxed{
v_R=
\frac{vP_{\mathrm{cr}}}
{RT_{\mathrm{cr}}}
}
\end{equation*}

Then calculate

\begin{equation*}
\boxed{
T_R=\frac{T}{T_{\mathrm{cr}}}
}
\end{equation*}

Using $v_R$ and $T_R$, the generalized compressibility chart can be used to determine the corresponding reduced pressure $P_R$.

The actual pressure is then obtained from

\begin{equation*}
\boxed{
P=P_RP_{\mathrm{cr}}
}
\end{equation*}

The compressibility factor does not necessarily have to be determined explicitly when the chart directly provides the required reduced property.

\textbf{\\ Accuracy of the Generalized Compressibility Chart}

The generalized compressibility chart provides approximate rather than exact values.

The accuracy depends on:
\begin{itemize}
    \item the resolution of the chart,
    \item the accuracy with which the chart is read, and
    \item how well the substance follows the principle of corresponding states.
\end{itemize}

\newpage

\textbf{Question: Why not define reduced specific volume as $v_R=v/v_{\mathrm{cr}}$ instead of using the pseudo-critical volume?}

Yes, we could define

\[
\boxed{v_R^*=\frac{v}{v_{\mathrm{cr}}}}
\]

There is nothing mathematically wrong with this definition. The reason the compressibility-factor approach instead uses

\[
\boxed{
v_R=\frac{vP_{\mathrm{cr}}}{RT_{\mathrm{cr}}}
}
\]

is that it makes the reduced properties fit naturally into the definition of the compressibility factor.

\textbf{\\ Start with the definition of $Z$}

\[
\boxed{
Z=\frac{Pv}{RT}
}
\]

The reduced pressure and temperature are defined as

\[
\boxed{
P_R=\frac{P}{P_{\mathrm{cr}}}
}
\qquad
\boxed{
T_R=\frac{T}{T_{\mathrm{cr}}}
}
\]

Therefore,

\[
P=P_RP_{\mathrm{cr}},
\qquad
T=T_RT_{\mathrm{cr}}
\]

Substituting into the definition of $Z$,

\[
Z=
\frac{(P_RP_{\mathrm{cr}})v}
{T_RT_{\mathrm{cr}}}
\]

Hence,

\[
Z=
\frac{P_R}{T_R}
\frac{vP_{\mathrm{cr}}}{RT_{\mathrm{cr}}}
\]

The last term is defined as the reduced specific volume:

\[
\boxed{
v_R=
\frac{vP_{\mathrm{cr}}}{RT_{\mathrm{cr}}}
}
\]

Therefore,

\[
\boxed{
Z=\frac{P_Rv_R}{T_R}
}
\]

This gives a clean relationship among the three reduced properties.

\textbf{\\ What happens if we use the actual critical volume?}

Suppose we instead define

\[
\boxed{
v_R^*=\frac{v}{v_{\mathrm{cr}}}
}
\]

Then

\[
v=v_R^*v_{\mathrm{cr}}
\]

Substituting into

\[
Z=\frac{Pv}{RT}
\]

and using

\[
P=P_RP_{\mathrm{cr}},
\qquad
T=T_RT_{\mathrm{cr}},
\]

gives

\[
Z=
\frac{P_RP_{\mathrm{cr}}v_R^*v_{\mathrm{cr}}}
{T_RT_{\mathrm{cr}}}
\]

Therefore,

\[
Z=
\frac{P_Rv_R^*}{T_R}
\left(
\frac{P_{\mathrm{cr}}v_{\mathrm{cr}}}
{RT_{\mathrm{cr}}}
\right)
\]

But

\[
\boxed{
Z_{\mathrm{cr}}
=
\frac{P_{\mathrm{cr}}v_{\mathrm{cr}}}
{RT_{\mathrm{cr}}}
}
\]

Thus,

\[
\boxed{
Z=
Z_{\mathrm{cr}}
\frac{P_Rv_R^*}{T_R}
}
\]

So using $v/v_{\mathrm{cr}}$ introduces an additional factor, $Z_{\mathrm{cr}}$.

\textbf{\\ Why is this undesirable?}

The purpose of reduced properties is to obtain a convenient dimensionless description of different gases:

\[
\boxed{
Z=f(P_R,T_R)
}
\]

Using

\[
\boxed{
v_R=\frac{vP_{\mathrm{cr}}}{RT_{\mathrm{cr}}}
}
\]

allows the compressibility factor to be written directly as

\[
\boxed{
Z=\frac{P_Rv_R}{T_R}
}
\]

without requiring the additional substance-dependent quantity $Z_{\mathrm{cr}}$.

If we use

\[
v_R^*=\frac{v}{v_{\mathrm{cr}}},
\]

then

\[
\boxed{
Z=
Z_{\mathrm{cr}}
\frac{P_Rv_R^*}{T_R}
}
\]

and $Z_{\mathrm{cr}}$ must also be known.

\textbf{\\ The deeper reason}

The ideal-gas equation is

\[
Pv=RT
\]

so the natural volume scale associated with the critical pressure and temperature is

\[
\boxed{
v_{\mathrm{ref}}
=
\frac{RT_{\mathrm{cr}}}{P_{\mathrm{cr}}}
}
\]

Therefore,

\[
\boxed{
v_R
=
\frac{v}{v_{\mathrm{ref}}}
=
\frac{vP_{\mathrm{cr}}}{RT_{\mathrm{cr}}}
}
\]

This volume is called the \textbf{pseudo-critical specific volume}.

It is important to distinguish:

\[
\boxed{
v_{\mathrm{cr}}
=
\text{actual physical critical specific volume}
}
\]

from

\[
\boxed{
v'_{\mathrm{cr}}
=
\frac{RT_{\mathrm{cr}}}{P_{\mathrm{cr}}}
=
\text{constructed pseudo-critical specific volume}
}
\]

In general,

\[
\boxed{
v'_{\mathrm{cr}}\neq v_{\mathrm{cr}}
}
\]

because

\[
Z_{\mathrm{cr}}
=
\frac{P_{\mathrm{cr}}v_{\mathrm{cr}}}
{RT_{\mathrm{cr}}}
\neq1
\]

for a real gas.

\textbf{\\ Final idea to remember}

\[
\boxed{
\frac{v}{v_{\mathrm{cr}}}
\text{ is possible, but it introduces }Z_{\mathrm{cr}}.
}
\]

\[
\boxed{
\frac{vP_{\mathrm{cr}}}{RT_{\mathrm{cr}}}
\text{ is used because it gives }
Z=\frac{P_Rv_R}{T_R}
\text{ directly.}
}
\]

Thus, the pseudo-critical volume is not being used because the actual critical volume cannot be measured. It is used because

\[
\boxed{
\text{it is the volume scale that is consistent with }
Pv=RT
\text{ and the reduced-property formulation.}
}
\]
















\mysection{3.8}{Other Equations of State}{3.8}

The ideal-gas equation of state is very simple and convenient to use, but its range of applicability is limited.

It is desirable to have equations of state that represent the $P$-$v$-$T$ behavior of substances accurately over a larger region.

Such equations of state are naturally more complicated than the ideal-gas equation.

Several equations of state have been proposed for representing real-gas behavior. The important equations discussed in this section are:
\begin{itemize}
    \item van der Waals equation of state,
    \item Beattie--Bridgeman equation of state,
    \item Benedict--Webb--Rubin equation of state, and
    \item virial equation of state.
\end{itemize}

These equations are applicable to the \textbf{gas phase} of substances and should not be used for liquids or liquid--vapor mixtures.

\textbf{\\ 1. van der Waals Equation of State}

The van der Waals equation of state was one of the earliest attempts to improve upon the ideal-gas equation of state.

It was developed by including two effects that are neglected in the ideal-gas model:
\begin{itemize}
    \item intermolecular attraction forces, and
    \item the volume occupied by the gas molecules themselves.
\end{itemize}

The van der Waals equation is

\begin{equation*}
\boxed{
\left(P+\frac{a}{v^2}\right)(v-b)=RT
}
\end{equation*}

where:
\begin{itemize}
    \item $a$ accounts for \textbf{intermolecular attraction forces},
    \item $b$ accounts for the \textbf{volume occupied by the gas molecules},
    \item $P$ is pressure,
    \item $v$ is specific volume,
    \item $T$ is absolute temperature, and
    \item $R$ is the specific gas constant.
\end{itemize}

\textbf{\\ Correction for Intermolecular Forces}

The term

\begin{equation*}
\frac{a}{v^2}
\end{equation*}

accounts for the attractive forces between molecules.

In the ideal-gas model, intermolecular attraction forces are neglected.

As pressure increases, the effects of intermolecular forces become increasingly important, and the ideal-gas model becomes less accurate.

\textbf{\\ Correction for Molecular Volume}

The quantity

\begin{equation*}
v-b
\end{equation*}

is used instead of $v$ in the ideal-gas relation.

The constant $b$ represents the volume occupied by the gas molecules per unit mass.

At low pressure, the molecular volume is generally a very small fraction of the total volume.

As pressure increases, the volume occupied by the molecules becomes an increasingly significant fraction of the total volume.

\textbf{\\ Determination of van der Waals Constants}

\insertfigure{0.4}{Figure : 13.png}

The constants $a$ and $b$ are determined using critical-point behavior.

At the critical point, the critical isotherm on a $P$-$v$ diagram has a \textbf{horizontal inflection point}.

Therefore, at the critical point,

\begin{equation*}
\boxed{
\left(\frac{\partial P}{\partial v}\right)_{T=T_{\mathrm{cr}}}=0
}
\end{equation*}

and

\begin{equation*}
\boxed{
\left(\frac{\partial^2P}{\partial v^2}\right)_{T=T_{\mathrm{cr}}}=0
}
\end{equation*}

Using these conditions, the van der Waals constants can be expressed in terms of the critical properties:

\begin{equation*}
\boxed{
a=\frac{27R^2T_{\mathrm{cr}}^2}{64P_{\mathrm{cr}}}
}
\end{equation*}

and

\begin{equation*}
\boxed{
b=\frac{RT_{\mathrm{cr}}}{8P_{\mathrm{cr}}}
}
\end{equation*}

Thus, the constants $a$ and $b$ can be determined from the critical-point data of a substance.

The van der Waals equation can also be expressed on a unit-mole basis by replacing:
\begin{itemize}
    \item $v$ with molar specific volume $\bar{v}$, and
    \item $R$ with the universal gas constant $R_u$.
\end{itemize}

\textbf{\\ Limitations of the van der Waals Equation}

The van der Waals equation has historical importance because it was one of the first equations to incorporate real-gas effects.

However, its accuracy is often inadequate.

Its accuracy can be improved by using values of $a$ and $b$ based on the actual behavior of the gas over a wider range rather than determining them from a single critical point.

Therefore, although the van der Waals equation provides a conceptual improvement over the ideal-gas equation, more sophisticated equations are generally preferred when higher accuracy is required.

\textbf{\\ 2. Beattie--Bridgeman Equation of State}

The Beattie--Bridgeman equation of state is a more complex equation designed to represent real-gas behavior more accurately.

It is based on experimentally determined constants.

The equation is

\begin{equation*}
\boxed{
P=
\frac{R_uT}{\bar{v}^{\,2}}
\left(1-\frac{c}{\bar{v}T^3}\right)
(\bar{v}+B)
-\frac{A}{\bar{v}^{\,2}}
}
\end{equation*}

where the quantities $A$ and $B$ are defined by

\begin{equation*}
\boxed{
A=A_0\left(1-\frac{a}{\bar{v}}\right)
}
\end{equation*}

and

\begin{equation*}
\boxed{
B=B_0\left(1-\frac{b}{\bar{v}}\right)
}
\end{equation*}

The constants appearing in the equation are experimentally determined for the substances of interest.

The required constants are tabulated for various substances.

\textbf{\\ Applicability of the Beattie--Bridgeman Equation}

The Beattie--Bridgeman equation provides reasonably accurate results over a useful range of gas densities.

Its applicability is generally expressed relative to the critical density:

\begin{equation*}
\boxed{
\rho\leq0.8\rho_{\mathrm{cr}}
}
\end{equation*}

where:
\begin{itemize}
    \item $\rho$ is the density of the gas, and
    \item $\rho_{\mathrm{cr}}$ is the density at the critical point.
\end{itemize}

The equation is therefore more accurate over a wider range than the van der Waals equation, but it still has limitations.

\textbf{\\ 3. Benedict--Webb--Rubin Equation of State}

The Benedict--Webb--Rubin equation of state was developed as an extension of the Beattie--Bridgeman equation.

It provides greater flexibility by using a larger number of experimentally determined constants.

The equation is

\begin{equation*}
\boxed{
\begin{aligned}
P={}&
\frac{R_uT}{\bar{v}}
+
\left(
B_0R_uT-A_0-\frac{C_0}{T^2}
\right)\frac{1}{\bar{v}^{\,2}}
\\
&+
\frac{bR_uT-a}{\bar{v}^{\,3}}
+
\frac{a\alpha}{\bar{v}^{\,6}}
+
\frac{c}{\bar{v}^{\,3}T^2}
\left(
1+\frac{\gamma}{\bar{v}^{\,2}}
\right)
e^{-\gamma/\bar{v}^{\,2}}
\end{aligned}
}
\end{equation*}

where the constants

\begin{equation*}
A_0,\quad B_0,\quad C_0,\quad a,\quad b,\quad c,\quad \alpha,\quad \gamma
\end{equation*}

are determined experimentally for the substance.

The values of these constants are tabulated for various gases.

\textbf{\\ Applicability of the Benedict--Webb--Rubin Equation}

The Benedict--Webb--Rubin equation can handle substantially higher gas densities than the Beattie--Bridgeman equation.

Its applicability is generally expressed as

\begin{equation*}
\boxed{
\rho\leq2.5\rho_{\mathrm{cr}}
}
\end{equation*}

Thus, the Benedict--Webb--Rubin equation can represent real-gas behavior over a relatively wide density range.

\textbf{\\ Strobridge Extension}

The Benedict--Webb--Rubin equation was subsequently extended by Strobridge.

The extension increases the number of constants and provides greater flexibility in representing the $P$-$v$-$T$ behavior of gases.

The resulting complex equations are particularly suitable for \textbf{digital computer calculations}.

\textbf{\\ 4. Virial Equation of State}

The equation of state of a substance can also be expressed as a series.

The general virial equation of state is

\begin{equation*}
\boxed{
P=
\frac{RT}{v}
+
\frac{a(T)}{v^2}
+
\frac{b(T)}{v^3}
+
\frac{c(T)}{v^4}
+
\frac{d(T)}{v^5}
+\cdots
}
\end{equation*}

The coefficients

\begin{equation*}
a(T),\quad b(T),\quad c(T),\quad d(T),\ldots
\end{equation*}

are called \textbf{virial coefficients}.

\textbf{\\ Virial Coefficients}

The virial coefficients are functions of temperature alone.

They can be determined:
\begin{itemize}
    \item experimentally, or
    \item theoretically using statistical mechanics.
\end{itemize}

Additional terms can be included in the virial equation to improve its representation of the $P$-$v$-$T$ behavior of a substance over a wider range.

Thus, the accuracy of the virial equation depends on the number of terms retained.

\textbf{\\ Low-Pressure Limit of the Virial Equation}

As pressure approaches zero, the virial coefficients vanish and the equation reduces to the ideal-gas equation of state.

Therefore,

\begin{equation*}
P\rightarrow0
\end{equation*}

leads to

\begin{equation*}
\boxed{
Pv\rightarrow RT
}
\end{equation*}

or equivalently,

\begin{equation*}
\boxed{
Z\rightarrow1
}
\end{equation*}

Thus, the virial equation is consistent with the ideal-gas limit.

\textbf{\\ Comparison of Equations of State}

The equations discussed in this section differ mainly in:
\begin{itemize}
    \item complexity,
    \item number of constants,
    \item range of applicability,
    \item accuracy, and
    \item suitability for computer calculations.
\end{itemize}

\begin{center}
\hspace*{-2cm}
\begin{tabular}{|c|c|c|}
\hline
\textbf{Equation} & \textbf{Main Characteristic} & \textbf{Relative Complexity} \\
\hline
van der Waals & Accounts for molecular attraction and molecular volume & Low \\
\hline
Beattie--Bridgeman & Uses experimentally determined constants & Moderate \\
\hline
Benedict--Webb--Rubin & More accurate and applicable over a wider range & High \\
\hline
Strobridge extension & Further increases the number of constants & Very high \\
\hline
Virial & Series representation with temperature-dependent coefficients & Variable \\
\hline
\end{tabular}
\end{center}

\textbf{\\ Accuracy and Complexity}

There is generally a trade-off between the simplicity and accuracy of an equation of state.

The van der Waals equation:
\begin{itemize}
    \item is simple,
    \item requires only two constants, and
    \item has a relatively limited range of accuracy.
\end{itemize}

The Beattie--Bridgeman equation:
\begin{itemize}
    \item is more complex,
    \item uses experimentally determined constants, and
    \item provides improved accuracy over an appropriate density range.
\end{itemize}

The Benedict--Webb--Rubin equation:
\begin{itemize}
    \item is more complex,
    \item uses more constants,
    \item represents the $P$-$v$-$T$ behavior more accurately, and
    \item is applicable over a wider density range.
\end{itemize}

The virial equation:
\begin{itemize}
    \item represents the equation of state as a series,
    \item uses temperature-dependent virial coefficients, and
    \item can be made more accurate by including additional terms.
\end{itemize}

\textbf{\\ Suitability for Hand Calculations and Computers}

Complex equations of state can represent the $P$-$v$-$T$ behavior of substances reasonably well and are particularly suitable for \textbf{digital computer applications}.

For hand calculations, however, property tables or simpler equations of state are generally more convenient.

This is especially important for specific-volume calculations because the equations discussed are generally \textbf{implicit in specific volume}.

Consequently, determining specific volume may require a \textbf{trial-and-error procedure}.


\textbf{\\ Important Limitation}

The equations of state discussed in this section are intended for the \textbf{gas phase}.

They should not be used for:

\begin{equation*}
\boxed{
\text{liquids}
}
\end{equation*}

or

\begin{equation*}
\boxed{
\text{liquid--vapor mixtures}
}
\end{equation*}

Property tables or other appropriate relations should be used for those regions.

\textbf{\\ General Comparison}

\begin{center}
\begin{tabular}{|c|c|c|c|}
\hline
\textbf{Equation} & \textbf{Constants} & \textbf{Accuracy} & \textbf{Use} \\
\hline
van der Waals & Few & Limited & Simple real-gas modeling \\
\hline
Beattie--Bridgeman & Several & Reasonable & Wider gas-density range \\
\hline
Benedict--Webb--Rubin & More & High & Accurate real-gas calculations \\
\hline
Strobridge & Many & High & Computer-oriented calculations \\
\hline
Virial & Variable & Depends on terms & Series-based representation \\
\hline
\end{tabular}
\end{center}

























\end{document}